% ============================================
% LUC-07
% Gerado de molambudos.md — build_miolo.py
% ============================================
\section*{\textbf{LUC-07} --- O Rastro de Oliveira}
\markboth{LUC-07 --- O Rastro de Oliveira}{}

\noindent \textit{Seguindo o ciclo reverso}

\noindent \textbf{Abril de 2026 --- Belo Horizonte / Quixeramobim}

\vspace{0.5em}\noindent\rule{\textwidth}{0.5pt}\vspace{0.5em}

\noindent Depois de ler a carta de Oliveira, Lúcia fez o que qualquer psiquiatra forense faria: seguiu o rastro.

\noindent O Dr. Heitor Oliveira desapareceu em 13 de junho de 1979. Seu carro foi encontrado três semanas depois num estacionamento em Quixeramobim, Ceará. A gasolina tinha acabado. As portas estavam destrancadas. Dentro, uma mochila com roupas, um caderno com 32 páginas escritas (que ninguém leu), e uma passagem de ônibus para Barbacena com data de 17 de junho --- o mesmo dia em que ele chegou.

\noindent Ele tinha voltado.

\noindent Lúcia passou uma tarde no Arquivo Público, revirando caixas. Encontrou o caderno de Oliveira --- esquecido, empoeirado, com a etiqueta \textquotedblleft{}Pertences do Dr. H. Oliveira --- 1979\textquotedblright{}.

\noindent Ela leu as 32 páginas em 40 minutos.

\noindent Eram anotações clínicas, mas com uma diferença: a letra mudava ao longo das páginas. Começava firme, médica. Terminava trêmula, irregular, \textit{infantil}.

\noindent A última frase era:

\noindent \textit{\textquotedblleft{}Voltei para Barbacena. O ciclo não me deixa ir. Devolvo o diário ao arquivo. Espero que a próxima leitora entenda melhor do que eu.\textquotedblright{}}

\noindent \textbf{A próxima leitora.}

\noindent Ela.

\noindent Ela era a próxima.

\noindent A fila tinha chegado nela.

\vspace{0.5em}\noindent\rule{\textwidth}{0.5pt}\vspace{0.5em}

\noindent Lúcia passou a noite no arquivo.

\noindent Leu o diário de Oliveira. Leu os relatórios das sessões. Leu a carta que Oliveira deixou para ela --- escrita em 1979, entregue em 2026.

\noindent E, pela segunda vez, leu o diário de Joaquim.

\noindent Só que, dessa vez, ela leu com outros olhos.

\noindent Não como psiquiatra.
Como \textit{herdeira}.

\noindent Cada palavra entrava. Cada frase ecoava. Ela sentiu o cheiro adocicado subir das páginas. O olho direito começou a arder.

\noindent \textit{\textquotedblleft{}O ciclo não tem volta\textquotedblright{}}, disse uma voz dentro dela.

\noindent Não era a voz dela.

\noindent Ela fechou o diário.

\noindent Guardou na mochila.

\noindent E, pela primeira vez, teve certeza de que não ia conseguir parar.

\vspace{0.5em}
\noindent\textit{\textbf{↪ Rotas:} \rota{LUC-08} \rotasep \rota{DOC-11}}